\documentclass[a4paper, 12pt]{exam}

\usepackage[onehalfspacing]{setspace}
\usepackage{graphicx}
\usepackage{amsmath,amsfonts,amsthm,amssymb,mathtools}
\usepackage{nameref}
\usepackage{hyperref}
\usepackage{xcolor}
\usepackage{gensymb}
\usepackage[]{booktabs}
\usepackage[utf8]{inputenc}
\usepackage{array}
\usepackage{setspace}
\usepackage{verbatim} %for multiline comments https://tex.stackexchange.com/questions/44282/multiline-comment
\usepackage{xhfill}
\usepackage{enumitem}
\usepackage{varwidth}
\usepackage{multicol, array}
\usepackage{etoolbox}
\usepackage{multicol}
\usepackage[most,many,breakable]{tcolorbox}
\usepackage{mathtools}
\usepackage{background} %Package for adding a watermark
\usepackage{tikz-cd}
\usepackage{pgfplots} 
\pgfplotsset{compat=newest}
\usepgfplotslibrary{patchplots}
\usepackage{anyfontsize}
\usepackage{sectsty}
\usepackage[space]{grffile} %To allow for spaces in the path for the watermark
\usepackage{amsmath}
\usepackage{geometry} %lets us set the margins of the page
\usepackage{comment} %Allows for multi-line comments (\ifx \fi)
\usepackage{import} %lets us perform multi-file compilation
\usepackage{pdfpages} %allows for the inclusion of external multi-page PDF documents
\usepackage{algpseudocode} %lets you write algorithms (probably won't be used but i'm including it anyways) https://www.overleaf.com/learn/latex/Algorithms
\usepackage{bookmark}
\usepackage{theorem}

\geometry{left=0.8in,top=0.5in, right = 0.8in}

\backgroundsetup{contents=\includegraphics{Anant_Logo.png}, opacity=0.1, angle=0, scale=1.5}

\begin{document}
	\title{Anant 24-25 Recruitment Test}
	\author{Attitude Determination and Controls Subsystem (ADCS)}
	\date{January 12th, 2024}
	\maketitle
	
	\section*{Instructions}
	All general instructions are provided in the Google Classroom. Make sure you refer to them for submission details. Failure to adhere to the given rules may result in disqualification. 
	\begin{enumerate}[]
		\item This paper relies more on comprehension and general aptitude than it does in depth knowledge of the subject. You are encouraged to use the internet to solve this paper.
		\item Any changes to the paper will be posted in Google Classroom and an updated file will be shared. Please keep an eye out.
	\end{enumerate}
	Rules of engagement
	\begin{enumerate}
		\item The goal of this paper is to test your general aptitude and comp
		\item 
		\item Collaboration with peers is prohibited. If you simply copy answers without understanding, you gain no knowledge unto yourself. Remember that you're writing this test for yourself, not for your friends. 
		\item You are free to break all the rules we've just listed here, but be warned that the interview process will be rigorous, and if you don't perform, you will be removed.
		\item If you have any questions regarding the content of the paper, feel free to reach out to any of the undersigned through WhatsApp for help.
		\item Contact Details:
		\begin{enumerate}
			\item Aaditya Thakkar - +91 81282 17114
			\item Nikhil Mathew - +91 97400 44411
			\item Pranav Chandra - +91 90807 96426
			\item Pramit Pal - +91 99026 97222
		\end{enumerate}
	\end{enumerate}
	\begin{center}
		All the best!
	\end{center}
	
	Our story revolves around Bob, a not-so-rebellious teen that's decided to try some green-rolly-rolly-puff-puff for the first time. Follow us and help him out as he drifts through space, trying not to get too \textit{high}.
	
	\pagebreak
	
	\newgeometry{left=1in,top=1.25in, right = 1in}

	\section*{Question 1 - Bob's Search for the Hill of Happiness}
	Bob, in his insanity of life, is looking for a mountain, Korok, which claims to have the best weed in the world. However, to find it, he must first decode a map, and achieve his state of high.
	
	\begin{enumerate}[label=(\alph*)]
		\item The map was created by a 4th dimensional being at 13:00 on the 2nd of January, 1920 in honor of a great man born on that day. He used what us mortals now call the Fixed Frame, which is defined as follows:
		\begin{center}
			\textit{The origin is defined to be the center of the Earth. The positive z-axis is the ray passing through the geographic north pole of the earth. The positive x-axis is in the direction of the vernal equinox, and the y-axis is taken as the right-hand cross product of the x and z axes.}
		\end{center}
		The coordinates in the map are $(0,132,155)$ in \textit{km}. Your job is to convert this into latitude and longitude so that Bob can figure out where on earth the mountain is. In your answers, include sketches of both the Fixed Frame and the Latitude Longitude frame, as well as a marking of points. (Assume each day is exactly 24 hours).
		
		\item Now, if you were unable to complete the first part of this question, that's alright. The higher beings have decided to give you the latitude-longitude-altitude to a Hill of Slightly-Less-Happiness, which Bob is okay with. However, Bob is way too far away from this point to get to on foot, so he'll be using his teleporter. However, his teleporter requires \textit{Rotating Frame} coordinates, defined as follows:
		\begin{center}
			\textit{The origin is the center of the Earth, and the x-axis is in from the center to the intersection between the prime meridian and the equator. The y-axis is the right-hand cross product of the x and z axes.}
		\end{center}
		Convert the latitude-longitude-altitude coordinates $(28°21'12.8"N 75°35'36.7"E)$ into the \textit{Rotating Earth} frame.
		
		\item 
	\end{enumerate}
	
	\pagebreak
	
	\section*{Question 2 - Bob's Fever Dream}
	Now that the weed has entered Bob's system, he's drifted into another plane. He's spiralling off of earth's surface and is now in an orbit with an altitude of $200km$. However, he wants to settle into an orbit at an altitude of $4206km$, using the least amount of energy. With his tiny rocket, he can only apply a burn of any impulse at three points in time. Find out the orbital transition and draw a diagram of where the burns should be made, and how much impulse should be provided. Hint:
	
	
	
	
	
	
\end{document}